\documentclass[10pt,a4paper,landscape]{article}
\usepackage[utf8]{inputenc}
\usepackage[T1]{fontenc}
\usepackage{lmodern}
\usepackage{xcolor}
\usepackage[margin=1.0cm]{geometry}
\usepackage{tikz}
\usetikzlibrary{positioning,arrows.meta,fit,calc}
\usepackage{parskip}
\pagestyle{empty}
\hyphenpenalty=3000 \emergencystretch=2em

\tikzset{
  blk/.style={draw, rounded corners=1.5pt, align=center, font=\footnotesize, inner sep=3pt, minimum height=8mm, fill=white},
  blkS/.style={blk, font=\scriptsize, minimum height=6mm},
  ext/.style={blk, dashed, fill=gray!8},
  grp/.style={draw, rounded corners=3pt, inner sep=7pt, font=\small\bfseries},
  moc/.style={-{Stealth}, very thick},
  sig/.style={-{Stealth}, thin},
  ster/.style={-{Stealth}, thin, dashed},
  bezp/.style={-{Stealth}, thick, red!70!black},
  lbl/.style={font=\tiny, fill=white, inner sep=1pt},
  krok/.style={draw, rounded corners=1.5pt, align=center, font=\scriptsize, inner sep=2.5pt, minimum height=7mm, fill=white, text width=44mm},
  bramka/.style={krok, fill=yellow!15},
  stopblk/.style={krok, fill=red!10, text width=46mm},
  kontakt/.style={draw, align=center, font=\scriptsize, inner sep=2.5pt, minimum height=7mm, fill=white}
}

\begin{document}

% =====================================================================
% STRONA TYTUŁOWA
% =====================================================================
\begin{center}
\vspace*{2.5cm}
{\LARGE\bfseries AMPERE POINT --- custom device}\\[6pt]
{\Large Schematy ideowe --- wersja 1}\\[14pt]
{\normalsize 3 lipca 2026}\\[22pt]
\begin{tabular}{ll}
\textbf{Arkusz 1} & Schemat blokowy całości (topologia Z1/Z2: przelot mocy z~odczepem sterującym, dwa moduły)\\[2pt]
\textbf{Arkusz 2} & Sekwencja pomiarowa z~bramkami i~ścieżkami błędów\\[2pt]
\textbf{Arkusz 3} & Łańcuch bezpieczeństwa (sprzętowy, niezależny od MCU)\\
\end{tabular}\\[20pt]
\begin{minipage}{0.72\textwidth}\small
Podstawa: \emph{AMPERE\_POINT\_custom\_device\_koncepcja\_v1.pdf} (zatwierdzona z~punktami Z1--Z6, 2026-07-03).
Schematy ideowe --- poziom architektury; wartości zabezpieczeń, przekroje i~oznaczenia aparatów zostaną
uszczegółowione w~schemacie elektrycznym (EPLAN) po testach etapu~0.\\[6pt]
\textbf{Legenda:} linia gruba $=$ tor mocy (L1/L2/L3/N/PE); linia cienka $=$ sygnał pomiarowy;
linia przerywana $=$ sterowanie / magistrale cyfrowe; linia czerwona $=$ łańcuch bezpieczeństwa /
ścieżka błędu; blok szary przerywany $=$ element zewnętrzny (DUT, przyrząd wzorcowany).
\end{minipage}
\end{center}
\newpage

% =====================================================================
% ARKUSZ 1 — SCHEMAT BLOKOWY
% =====================================================================
\noindent{\small\textbf{Arkusz 1/3 --- Schemat blokowy całości} \hfill custom device v1 --- 2026-07-03}\par\vspace{0.5mm}
\begin{center}
\begin{tikzpicture}[x=1cm,y=1cm]

% ---------- zewnętrzne (lewa kolumna) ----------
\node[ext, text width=26mm] (inst) at (1.6,17.3) {Instalacja zasilająca obiektu (TN/TT)};
\node[ext, text width=26mm, minimum height=11mm] (dut) at (1.6,15.1) {\textbf{Ładowarka badana (DUT)}\\ P/B 1F 32 A\\ Q11 3F 16 A};
\node[ext, text width=24mm] (ut595) at (1.6,11.0) {UT595 --- izolacja, ciągłość, Zs, RCD A (wzorcowany)};
\node[ext, text width=24mm] (ut522) at (1.6,5.6) {UT522 --- uziom 3-elektrodowy};
\draw[moc] (inst) -- (dut);

% ---------- tor mocy (górny rząd) ----------
\node[blk, text width=28mm] (inlet) at (5.9,16.2) {Gniazdo pojazdowe Type 2 32 A 3F\\ + blokada wtyku};
\node[blkS] (q0) at (8.9,16.2) {Rozłącznik\\ Q0};
\node[blkS] (fuse) at (11.35,16.2) {Bezpieczniki\\ gG 32 A};
\node[blkS, text width=19mm] (cdsr) at (13.75,16.2) {CDSR --- RCM\\ (okno L1L2L3N)};
\node[blkS] (ct) at (16.1,16.2) {3$\times$CT\\ 40 A};
\node[blkS] (k1) at (18.4,16.2) {Stycznik\\ główny K1};
\node[blkS, text width=18mm] (cee) at (21.1,16.2) {Gniazdo CEE\\ 32 A 5p};

\draw[moc] (dut.east) -- ++(0.45,0) |- node[lbl, pos=0.75, above]{kabel Type 2 DUT} (inlet.west);
\draw[moc] (inlet) -- (q0);
\draw[moc] (q0) -- (fuse);
\draw[moc] (fuse) -- (cdsr);
\draw[moc] (cdsr) -- (ct);
\draw[moc] (ct) -- (k1);
\draw[moc] (k1) -- (cee);

% ---------- PM701E + serwa + bananowe ----------
\node[blkS, text width=32mm] (serwa) at (5.0,14.25) {2$\times$ serwo DS3218 + rama (PWM)};
\node[blk, text width=38mm, minimum height=12mm] (pm) at (7.4,12.7) {\textbf{PM701E} (nienaruszony)\\ pokrętła: stan BB/C/A/B/D,\\ prąd NC/13/20/32/63 A};
\node[blkS, text width=28mm] (banan) at (7.4,10.8) {gniazda bananowe\\ L1/N/PE/L2/L3};
\node[blkS, text width=18mm] (esh) at (4.0,6.9) {zaciski\\ E/S/H};

\draw[sig] (inlet.south) -- ++(0,-0.45) node[lbl, right, yshift=-1mm]{odczep Type 2 (CP/PP/PE/L/N, $\le$1 m)} -| (pm.north);
\draw[ster] (serwa.east) -| node[lbl, pos=0.9, right]{na pokrętłach} ($(pm.north west)+(0.5,0)$);
\draw[sig] (pm) -- (banan);
\draw[sig, dashed] (banan.west) -- ++(-1.0,0) |- (ut595.east);
\draw[sig, dashed] (esh.west) -- ++(-0.5,0) |- (ut522.east);

% ---------- kolumny analogowe ----------
\node[blkS, text width=28mm] (separ) at (11.9,13.9) {przekaźniki separujące\\ (tryb izolacji)};
\node[blkS, text width=28mm] (gen) at (11.9,12.3) {generator upływu\\ DC 0--20 mA / AC (L$\rightarrow$PE)\\ bocznik + AMC1301};
\node[blkS, text width=28mm] (cp) at (11.9,10.7) {akwizycja CP/PP\\ (dzielnik, komparator)};
\node[blkS, text width=28mm] (met) at (15.8,13.9) {metering 3F --- SPI\\ ATM90E36A + dzielniki};
\node[blkS, text width=28mm] (irx) at (15.8,12.3) {termowizja --- I2C\\ MLX90640 + 2$\times$MLX90614};
\node[blkS, text width=28mm] (temp) at (15.8,10.7) {DS18B20 --- 1-Wire\\ szyny, K1, otoczenie};

\draw[sig] (fuse.south) -- (fuse.south|-separ.north);
\draw[sig] (separ) -- (gen);
\draw[sig] (ct.south) -- node[lbl, right, pos=0.5]{prądy} (ct.south|-met.north);
\draw[sig] ($(pm.east)+(0,-0.4)$) -- (9.65,12.3) -- node[lbl, left]{SIGNAL OUTPUT (CP)} (9.65,10.7) -- (cp.west);
\draw[ster] (cp.south) -- node[lbl, right, pos=0.3]{poziomy + duty} (cp.south|-11.9,8.45);

% ---------- magistrala cyfrowa ----------
\draw[ster] (met.east) -- ++(0.4,0) coordinate (busTR);
\draw[ster] (irx.east) -- ++(0.4,0);
\draw[ster] (temp.east) -- ++(0.4,0);
\draw[ster] (cdsr.south) -- node[lbl, left, pos=0.15]{SPI (RCM)} (13.75,8.85);
\draw[ster] (busTR) -- (17.7,8.85) -- node[lbl, above, pos=0.5]{SPI / I2C / 1-Wire} (13.4,8.85) -- (13.4,8.45);

% ---------- logika ----------
\node[blk, text width=38mm, fill=red!5, draw=red!70!black] (chain) at (6.3,7.7) {\textbf{łańcuch bezpieczeństwa 24 V}\\ (arkusz 3) $\rightarrow$ cewka K1};
\node[blk, text width=40mm, minimum height=13mm] (esp) at (11.6,7.7) {\textbf{ESP32-S3} --- sekwencer\\ + MCP23017 + ULN2803 (opto)\\ + microSD + RTC};
\node[blkS, text width=28mm] (hmi) at (16.4,8.3) {HMI DWIN 7$''$ --- UART\\ START / E-STOP / LED};
\node[blkS, text width=28mm] (sd) at (16.4,7.1) {protokoły CSV/JSON\\ (SD, znaczniki RTC)};
\node[blkS, text width=28mm] (wifi) at (16.4,5.9) {WiFi / MQTT\\ $\rightarrow$ Home Assistant};
\node[blkS, text width=28mm] (psu) at (20.8,8.3) {zasilacz 24/5/3,3 V\\ $\leftarrow$ przewód Schuko\\ (logika nie z~DUT)};

\draw[ster] (esp.east|-hmi.west) -- (hmi.west);
\draw[ster] (esp.east|-sd.west) -- (sd.west);
\draw[ster] ($(esp.south)+(0.9,0)$) |- (wifi.west);
\draw[sig] (psu.south) -- ++(0,-0.45) node[lbl, below]{zasilanie logiki i~cewek};
\draw[ster] (esp.west) -- node[lbl, above]{enable + watchdog} (chain.east);
\draw[bezp] (chain.north) -- (6.3,9.15) -- node[lbl, above, pos=0.18]{zezwolenie $\rightarrow$ cewka K1} (18.4,9.15) -- (k1.south);
\draw[ster] ($(esp.north)+(-1.8,0)$) -- (9.8,9.45) -- node[lbl, above, pos=0.35]{PWM serw} (4.6,9.45) -- (4.6,13.9);

% ---------- moduł B ----------
\node[blkS, text width=30mm] (wtyk) at (25.0,16.2) {wtyk CEE 32 A 5p};
\node[blk, text width=30mm] (stycz) at (25.0,14.5) {12$\times$ stycznik 25 A\\ (4 stopnie/fazę, cewki 24 V)};
\node[blk, text width=30mm, minimum height=11mm] (grzalki) at (25.0,12.5) {grzałki żebrowane L1/L2/L3\\ 1,0/1,5/2,0/3,0 kW\\ ($\sim$7,5 kW/fazę)};
\node[blkS, text width=30mm] (ksd) at (25.0,10.6) {KSD301 na sekcjach\\ (NC, szeregowo)};
\node[blkS, text width=30mm] (fan) at (25.0,9.3) {2$\times$ wentylator 500 m$^3$/h\\ + potwierdzenie pracy};
\node[blkS, text width=30mm] (harting) at (25.0,8.0) {złącze sterownicze\\ (cewki, termiki, wentyl.)};

\draw[moc] (cee.east) -- node[lbl, below=1.5mm]{5$\times$6 mm$^2$} (wtyk.west);
\draw[moc] (wtyk) -- (stycz);
\draw[moc] (stycz) -- (grzalki);
\draw[bezp] (ksd) -- (fan);
\draw[bezp] (fan) -- (harting);
\draw[bezp] (harting.west) -- ++(-0.25,0) -- (23.25,4.4) -- node[lbl, above, pos=0.55]{przewód sterowniczy: termiki + potwierdzenie wentylatorów $\rightarrow$ łańcuch (ark. 3)} (6.3,4.4) -- (chain.south);
\draw[ster] (esp.south) -- (11.6,4.9) -- node[lbl, above, pos=0.68]{24 V --- cewki stopni obciążenia} (22.9,4.9) -- (22.9,14.5) -- (stycz.west);

% ---------- kontenery ----------
\node[grp, fit=(inlet)(cee)(serwa)(pm)(banan)(esh)(separ)(gen)(cp)(met)(irx)(temp)(chain)(esp)(hmi)(sd)(wifi)(psu)] (modA) {};
\node[grp, fit=(wtyk)(stycz)(grzalki)(ksd)(fan)(harting)] (modB) {};
\node[font=\small\bfseries, anchor=south west] at ($(modA.north west)+(0.15,0.03)$) {MODUŁ A --- jednostka sterująco-pomiarowa (IP54)};
\node[font=\small\bfseries, anchor=south east] at ($(modB.north east)+(-0.1,0.03)$) {MODUŁ B --- obciążenie};

\end{tikzpicture}
\end{center}
\newpage

% =====================================================================
% ARKUSZ 2 — SEKWENCJA POMIAROWA
% =====================================================================
\noindent{\small\textbf{Arkusz 2/3 --- Sekwencja pomiarowa (bramki i~ścieżki błędów)} \hfill custom device v1 --- 2026-07-03}\par\vspace{0.5mm}
\begin{center}
\begin{tikzpicture}[x=1cm,y=1cm, node distance=3.5mm]

% kolumna 1
\node[krok, fill=green!12] (s0) at (3.1,15.8) {\textbf{0. SELF-TEST}\\ łańcuch bezp., wentylatory, czujniki,\\ kalibracja serw (BB/NC), komunikacja};
\node[bramka, below=of s0] (s1) {\textbf{1. TRYB IZOLACJI} --- separacja własnej\\ elektroniki i~modułu B od magistrali;\\ UT595 przez PM701E: izolacja 250 V\\ \textbf{bramka: $\ge$1 M$\Omega$ (wpis HMI)}};
\node[bramka, below=of s1] (s2) {\textbf{2. CIĄGŁOŚĆ PE} $\ge$200 mA\\ (UT595, wyjście PE)\\ \textbf{bramka: wpis wyniku}};
\node[krok, below=of s2] (s3) {\textbf{3. STAN B} (serwo stanu $\rightarrow$ B)\\ potwierdzenie CP; pre-testy PM701E};
\node[bramka, below=of s3] (s4) {\textbf{4. STAN C} (serwo $\rightarrow$ C, CP OK)\\ stycznik DUT zamyka; weryfikacja napięć\\ i~kolejności faz; blokada wtyku};

% kolumna 2
\node[krok] (s5) at (10.3,15.8) {\textbf{5a. RCD typ A protokołowo} (UT595):\\ I$_{\Delta n}$/2, I$_{\Delta n}$, 5I$_{\Delta n}$, czasy, U$_B$; wpis};
\node[krok, below=of s5] (s5b) {\textbf{5b. RDC-DD 6 mA DC auto}:\\ rampa generatora, rejestracja prądu\\ i~czasu zadziałania (zanik L)};
\node[krok, below=of s5b] (s5c) {po każdym zadziałaniu: prowadzony\\ \textbf{reset DUT} $\rightarrow$ powrót do stanu C};
\node[krok, below=of s5c] (s6) {\textbf{6. PĘTLA Zs} (UT595, gniazda\\ bananowe); warunek Zs$\cdot$Ia $\le$ U$_0$; wpis};
\node[bramka, below=of s6] (s7a) {\textbf{7. OBCIĄŻENIE}: serwo prądu $\rightarrow$\\ deklaracja (32 A 1F / 16 A 3F);\\ \textbf{bramka: łańcuch zamknięty\\ + START operatora} $\rightarrow$ K1};

% kolumna 3
\node[krok] (s7b) at (17.3,15.8) {stopnie 25/50/75/100\% deklaracji\\ (kombinacje $\le$ prąd z~CP)};
\node[krok, below=of s7b] (s7c) {\textbf{SOAK 15--30 min}: log U/I/P per faza,\\ asymetria, $\Delta T$ wtyku (MLX),\\ temperatury wewn.; kryteria pass/fail};
\node[krok, below=of s7c] (s8) {\textbf{8. TESTY BŁĘDÓW}: CP Error „E'',\\ PE Error (PM701E, etap 1 ręcznie);\\ pomiar czasu odcięcia DUT};
\node[krok, below=of s8] (s9) {\textbf{9. UZIOM} (niezależnie czasowo):\\ UT522 na E/S/H; wpis do protokołu};
\node[krok, fill=green!12, below=of s9] (s10) {\textbf{10. RAPORT}: CSV/JSON na SD,\\ MQTT $\rightarrow$ HA; protokół PDF na PC};

% przepływ główny
\draw[sig] (s0) -- (s1);
\draw[sig] (s1) -- node[lbl, right]{T} (s2);
\draw[sig] (s2) -- node[lbl, right]{T} (s3);
\draw[sig] (s3) -- (s4);
\draw[sig] (s4.south) -- ++(0,-0.3) -| (6.8,15.8) -- (s5.west);
\draw[sig] (s5) -- (s5b);
\draw[sig] (s5b) -- (s5c);
\draw[sig] (s5c) -- (s6);
\draw[sig] (s6) -- (s7a);
\draw[sig] (s7a.south) -- ++(0,-0.3) -| (13.9,15.8) -- (s7b.west);
\draw[sig] (s7b) -- (s7c);
\draw[sig] (s7c) -- (s8);
\draw[sig] (s8) -- (s9);
\draw[sig] (s9) -- (s10);

% blok przerwania + ścieżki błędów
\node[stopblk, minimum height=16mm] (stop) at (23.7,10.2) {\textbf{PRZERWANIE / STAN BEZPIECZNY}\\ rozpad zezwolenia $\rightarrow$ K1 otwarty,\\ zrzut obciążenia, serwa $\rightarrow$ BB/NC,\\ komunikat HMI, log przyczyny;\\ wznowienie tylko przez START};

\draw[bezp] (s1.west) -- node[lbl, above, pos=0.4]{N} (s1.west -| 0.5,0) -- (0.5,5.3) -- (23.7,5.3) -- (stop.south);
\draw[bezp] (s2.west) -- node[lbl, above, pos=0.4]{N} (s2.west -| 0.5,0);
\draw[bezp] (s4.west) -- node[lbl, above, pos=0.4]{CP$\neq$komenda} (s4.west -| 0.5,0);
\draw[bezp] (s7c.east) -- ++(0.4,0) -| node[lbl, above, pos=0.25]{anomalia} ($(stop.north)+(-0.9,0)$);
\draw[bezp] (s8.east) -- ++(0.7,0) -| node[lbl, above, pos=0.25]{brak reakcji DUT} ($(stop.north)+(0.9,0)$);

\node[blkS, text width=44mm, align=left] at (17.3,6.8) {\scriptsize\textbf{Warunki ciągłe (każdy krok):} CP zgodne z~komendą; RCM (CDSR) poniżej progu; termiki OK; wentylatory OK; E-STOP niewciśnięty. Naruszenie $\rightarrow$ PRZERWANIE.};

\end{tikzpicture}
\end{center}
\newpage

% =====================================================================
% ARKUSZ 3 — ŁAŃCUCH BEZPIECZEŃSTWA
% =====================================================================
\noindent{\small\textbf{Arkusz 3/3 --- Łańcuch bezpieczeństwa (sprzętowy, niezależny od MCU)} \hfill custom device v1 --- 2026-07-03}\par\vspace{4mm}
\begin{center}
\begin{tikzpicture}[x=1cm,y=1cm]

% ---------- szczebel 1: łańcuch zezwolenia -> cewka RB ----------
\node[font=\small\bfseries] at (0.7,13.4) {+24 V};
\node[kontakt, text width=20mm] (estop) at (3.0,13.4) {\textbf{E-STOP}\\ grzybek, NC (panel A)};
\node[kontakt, text width=26mm] (ksd) at (6.3,13.4) {\textbf{KSD301 $\times$12} --- NC\\ szeregowo (moduł B)};
\node[kontakt, text width=22mm] (lid) at (9.6,13.4) {\textbf{krańcówka} pokrywy\\ strefy mocy (NC)};
\node[kontakt, text width=26mm] (fans) at (12.9,13.4) {\textbf{potwierdzenie wentylatorów}\\ (przekaźnik prądowy)};
\node[kontakt, text width=24mm] (mcuen) at (16.2,13.4) {\textbf{zezwolenie MCU}\\ (może tylko rozłączyć)};
\node[kontakt, text width=20mm] (wdt) at (19.3,13.4) {\textbf{watchdog}\\ człon impulsowy};
\node[kontakt, text width=22mm] (start) at (22.4,14.4) {\textbf{START}\\ przycisk NO};
\node[kontakt, text width=22mm] (rbaux) at (22.4,12.4) {styk pomocniczy RB\\ (samopodtrzymanie)};
\node[kontakt, text width=20mm] (cewkarb) at (25.6,13.4) {\textbf{cewka przekaźnika}\\ \textbf{bezpieczeństwa RB}};

\draw[bezp] (1.1,13.4) -- (estop);
\draw[bezp] (estop) -- (ksd);
\draw[bezp] (ksd) -- (lid);
\draw[bezp] (lid) -- (fans);
\draw[bezp] (fans) -- (mcuen);
\draw[bezp] (mcuen) -- (wdt);
\draw[bezp] (wdt.east) -- (20.9,13.4);
\draw[bezp] (20.9,13.4) |- (start.west);
\draw[bezp] (20.9,13.4) |- (rbaux.west);
\draw[bezp] (start.east) -| (24.2,13.4);
\draw[bezp] (rbaux.east) -| (24.2,13.4);
\draw[bezp] (24.2,13.4) -- (cewkarb.west);
\draw[bezp] (cewkarb.east) -- (27.3,13.4) node[font=\small\bfseries, right]{0 V};

\node[lbl, text width=56mm, align=left] at (23.3,11.2) {\scriptsize samopodtrzymanie: po każdym rozpadzie łańcucha ponowne załączenie wymaga fizycznego wciśnięcia START};

% ---------- szczebel 2: styk główny RB -> cewka K1 ----------
\node[font=\small\bfseries] at (0.7,5.6) {+24 V};
\node[kontakt, text width=28mm] (rbmain) at (12.2,5.6) {\textbf{styk główny RB}\\ (wymuszony mechanicznie)};
\node[kontakt, text width=28mm] (cewkak1) at (16.6,5.6) {\textbf{cewka stycznika głównego K1}\\ (tor mocy, arkusz 1)};
\draw[bezp] (1.1,5.6) -- (rbmain.west);
\draw[bezp] (rbmain) -- (cewkak1);
\draw[bezp] (cewkak1.east) -- (19.4,5.6) node[font=\small\bfseries, right]{0 V};

% sprzezenie mechaniczne RB
\draw[ster] (cewkarb.south) -- (25.6,6.6) -- node[lbl, above, pos=0.5]{sprzężenie mechaniczne styków RB (styki wymuszone)} (12.2,6.6) -- (rbmain.north);

% ---------- MCU: odczyt + enable + watchdog ----------
\node[blkS, text width=56mm, align=left] (mcu) at (6.5,7.9) {\textbf{ESP32-S3 --- tylko odczyt i~zezwolenie:}\\ stan każdego członu przez optoizolację (diagnostyka na HMI); wyjście enable; generator impulsów watchdog (zanik impulsów $=$ rozpad); \emph{firmware nie może zmostkować żadnego elementu łańcucha}};

\draw[ster] (estop.south) -- (3.0,12.35) -- node[lbl, above, pos=0.35]{odczyt stanu członów (optoizolacja)} (9.0,12.35) -- (9.0,8.95);
\draw[ster] (ksd.south) -- (6.3,12.35);
\draw[ster] (lid.south) -- (9.6,12.35);
\draw[ster] (mcu.east) -- ++(0.4,0) -- (9.9,8.35) -- node[lbl, above, pos=0.6]{enable} (16.2,8.35) -- (mcuen.south);
\draw[ster] ($(mcu.east)+(0,-0.5)$) -- ++(0.55,0) -- (10.05,7.4) -- node[lbl, above, pos=0.6]{impulsy watchdog} (19.3,7.4) -- (wdt.south);

% ---------- uzupelnienia ----------
\node[blkS, text width=62mm, align=left] (uzup) at (5.0,3.4) {\textbf{Uzupełniająco (poza łańcuchem):}\\
$\bullet$ CDSR jako RCM --- próg awaryjny $\rightarrow$ zdjęcie enable\\
$\bullet$ bezpieczniki gG na fazach i~stopniach obciążenia\\
$\bullet$ rozłącznik główny Q0 (odcięcie ręczne)\\
$\bullet$ blokada elektromagnesowa wtyku Type 2 w~stanie C\\
$\bullet$ praca impulsowa/duty-cycle load banku (termika)};

\node[blkS, text width=58mm, align=left] at (13.6,3.4) {\textbf{Zasada:} tor 24 V cewki K1 przechodzi wyłącznie przez styki sprzętowe. Awaria dowolnego członu, zanik zasilania, zawieszenie firmware lub E-STOP $\Rightarrow$ RB odpada $\Rightarrow$ K1 otwiera tor mocy do modułu obciążenia.};

\end{tikzpicture}
\end{center}

\end{document}
